\cleardoubleoddpage%  Make sure to start each chapter on a new odd page
\chapter{Background}
\label{sec:Background}


\section{Reinforcement Learning}
RL is a branch of machine learning where an agent learns to make decisions by interacting with an environment using reward signals. At each step, the agent observes the current state, selects an action, and receives a reward as feedback. The objective of the agent is to learn a policy that maximizes the expected cumulative reward over time.

Most RL problems are modeled as Markov Decision Processes (MDPs) $(\mathcal{S}, \mathcal{A}, p, r, \gamma)$. Here, $\mathcal{S}$ denotes state space of the environment, $\mathcal{A}$ represents the action space, $p(s'|s,a)$ describes the probability of moving to state $s'$ after taking action $a$ in state $s$, $r(s,a)$ denotes the reward function that specifies the immediate reward received after taking action $a$ in state $s$, and $\gamma \in [0,1]$ is the discount factor that determines how future rewards are weighted relative to immediate rewards. The future accumulated reward, also known as return $R_t$ or gain $G_t$ is defined as:

$$R_t = \sum_{i= t+1}^\infty \gamma^i r(s_i, a_i, s_{i+1})$$

During training, the agent improves its policy by exploring the environment, exploiting previously collected knowledge and updating its estimates of action values based on the rewards it receives.




\section{Deep Reinforcement Learning}
Deep RL combines RL with deep NNs. In this setting, NNs are used as function approximators to estimate value functions or policies directly from high-dimensional inputs, allowing us to scale algorithms to more complex problems.

 Examples of Deep RL algorithms include Deep Q-Networks (DQN) \cite{DQN}, which approximate action-value functions, and actor-critic methods that learn both value functions and policies simultaneously. A number of algorithms like DQN learn using a temporal difference (TD)-based learning where we minimize the loss defined as follows:

 $$ \mathcal{L}(\theta) = l_k(r +\gamma \max_{a^`} Q_{\theta^`}(s^`, a^`) - Q_{\theta}(s,a))$$

Where $\theta$ refers to the parameter of the Q-network. In DQN the Huber loss is used defined as:

$$ l_{\kappa}(\delta) =
\begin{cases}
\frac{1}{2}\delta^2 & \text{if } |\delta| \le \kappa \\
\kappa\left(|\delta| - \frac{1}{2}\kappa\right) & \text{otherwise}
\end{cases}$$

When training any DQN-based variations of agents we use the Huber loss as defined above. However, it should be noted that all experiments described in this work are reproducible with other similar losses including but not limited to mean-square error (MSE) loss.

\section{Offline Reinforcement Learning}

Offline, or batch, RL \cite{lange2012batch} is a setting where the agent learns from a fixed dataset of previously collected state-action pairs instead of interacting directly with the environment during training. This dataset typically consists of tuples containing states, actions, rewards, and next states generated by a behavior policy.

Offline RL is particularly useful in domains where collecting new data is costly, dangerous, or impractical, such as healthcare, robotics, or autonomous driving. However, the absence of environment interaction introduces additional challenges like the extrapolation error. This phenomenon occurs when the agent overestimates values to out-of-distribution actions that were not observed in the dataset.



